\subsection{有限生成模的对偶性}

在本节中，我们考虑一个具有对偶化模 \(\Omega\) 的有限维 Noether 环 A。于是有 \(\mathrm{di}_{\Lambda}(\Omega) = \dim(\mathrm{A}) < +\infty\)（\(n^{\circ}\) 1, 命题 1）。选取一个有限长度的内射分解 \(e: \Omega \to (\mathbf{I}, \delta)\)。对 A-模的每个复形 C，用 \(\mathbf{D}(\mathbf{C})\) 表示复形 \(\operatorname{Homgr}_{\Lambda}(\mathbf{C}, \mathbf{I})\)。这特别适用于每个 A-模 M，把它视为集中在 0 次的复形；于是有 \(\mathbf{D}(\mathrm{M})^{i} = \mathrm{Hom}_{\mathrm{A}}(\mathrm{M},\mathrm{I}^{i})\)（对每个整数 \(i\)）。回忆在 A, X, 第 100 页，定理 1 中，我们构造了一个典范同构

\[
\varphi (\mathrm{M}, \mathrm{I}): \mathrm{H} (\mathbf {D} (\mathrm{M})) \longrightarrow \operatorname{Ext} _ {\mathrm{A}} (\mathrm{M}, \Omega).
\]

例.—— 1) 复形 \(\mathbf{D}(\mathbf{A}) = \operatorname{Homgr}_{\mathbf{A}}(\mathbf{A},\mathbf{l})\) 等同于 I。映射 \(c:\Omega \to \mathbf{D}(\Lambda)\) 按定义是一个拟同构。

\begin{enumerate}
\def\labelenumi{\arabic{enumi})}
\setcounter{enumi}{1}
\item
  同态 \(e \in \text{Homgr}_A(\Omega, I)^0\) 是 \(\mathbf{D}(\Omega)^0\) 的一个元素；A-线性映射 \(\tilde{e}: A \to \mathbf{D}(\Omega)\)（满足 \(\tilde{e}(1) = e\)）是一个拟同构（第 4 节，命题 7, a) 与 b)）。
\item
  设 S 为 A 的一个乘性子集。\(S^{-1}A\)-模 \(S^{-1}\Omega\) 是对偶化的（第 1 节，命题 2 的推论 1）；\(S^{-1}A\)-模 \(S^{-1}I^{i}\) 都是内射的（§ 3, 第 2 节，命题 3 的推论 1），而由 e 导出的态射 \(e^{\prime}:S^{-1}\Omega\to S^{-1}I\) 是 \(S^{-1}\Omega\) 的一个内射分解，因此可以对它应用前面的结果。对每个有限型的复形 C（特别地，对每个有限生成的 A-模 M），从 \(S^{-1}\mathbf{D}(\mathbf{C})\) 到 \(\operatorname{Homgr}_{\mathrm{S}^{-1}\mathrm{A}}(\mathrm{S}^{-1}\mathrm{C},\mathrm{S}^{-1}\mathrm{I})=\mathbf{D}(\mathrm{S}^{-1}\mathrm{C})\) 的典范同态是双射的。
\item
  设 A 为 Dedekind 环，K 为它的分式域。A-模 A 是对偶化的，且具有由正合列定义的长度为 1 的内射分解 I
\end{enumerate}

\[
0 \to \mathrm{A} \stackrel {e} {\longrightarrow} \mathrm{K} \stackrel {\delta} {\longrightarrow} \mathrm{K/A} \to 0
\]

其中 δ 是典范满射。对每个 A-模 M，复形 D(M) 是集中在次数 0 与 1 的复形

\[
\dots \longrightarrow 0 \longrightarrow \operatorname{Hom} _ {\mathrm{A}} (\mathrm{M}, \mathrm{K}) \stackrel {d} {\longrightarrow} \operatorname{Hom} _ {\mathrm{A}} (\mathrm{M}, \mathrm{K} / \mathrm{A}) \longrightarrow 0 \longrightarrow \dots
\]

其中 \(d = \mathrm{Hom}_{\mathrm{A}}(1_{\mathrm{M}},\delta)\)。我们有一个正合列

\[
0 \to \operatorname{Hom} _ {\mathrm{A}} (\mathrm{M}, \mathrm{A}) \longrightarrow \mathbf {D} (\mathrm{M}) ^ {0} \stackrel {d} {\longrightarrow} \mathbf {D} (\mathrm{M}) ^ {1} \longrightarrow \operatorname{Ext} _ {\mathrm{A}} ^ {1} (\mathrm{M}, \mathrm{A}) \to 0.
\]

对复形的每个态射 \(f: \mathrm{C} \to \mathrm{C}'\)，用 \(\mathbf{D}(f): \mathbf{D}(\mathrm{C}') \to \mathbf{D}(\mathrm{C})\) 表示复形的态射 \(\operatorname{Homgr}_{\mathrm{A}}(f, 1_{\mathrm{I}})\)。若 \(f\) 是拟同构，则 \(\mathbf{D}(f)\) 是拟同构（A, X, 第 86 页，命题 4, b)）。若 \(\mathrm{C}' \xrightarrow{f} \mathrm{C} \xrightarrow{g} \mathrm{C}''\) 是复形的一个正合列，则复形列 \(\mathbf{D}(\mathrm{C}''') \xrightarrow{\mathbf{D}(g)} \mathbf{D}(\mathrm{C}) \xrightarrow{\mathbf{D}(f)} \mathbf{D}(\mathrm{C}')\) 是正合的（A, X, 第 83 页，命题 2, a)）。

设 M 为 A-模。对 M 的每个元素 m，结合一个从 D(M) 到 I 的映射 \(\alpha_{\mathrm{M}}(m): f \mapsto f(m)\)；它是 \(\mathbf{D}(\mathbf{D}(\mathbf{M}))_{0} = \operatorname{Homgr}_{\mathrm{A}}(\mathbf{D}(\mathbf{M}), \mathrm{I})_{0}\) 的一个元素。由定义可知，\(\alpha_{\mathrm{M}}(m)\) 是复形的一个态射，从而是 \(\mathrm{Z}_{0}(\mathbf{D}(\mathbf{D}(\mathbf{M})))\) 的一个元素。

这样就定义了复形的一个态射：

\[
\alpha_ {\mathrm{M}}: \mathrm{M} \rightarrow \mathbf {D} (\mathbf {D} (\mathrm{M})) ,
\]

由此，通过取同调，得到 A-模的一个同态

\[
\alpha_ {\mathrm{M}}: \mathrm{M} \to \mathrm{H} _ {0} (\mathbf {D} (\mathbf {D} (\mathrm{M}))).
\]

定理 1. 设 M 为有限生成的 A-模。则 \(\alpha_{\mathrm{M}}\) 是一个拟同构：有 \(\mathrm{H}_{i}(\mathbf{D}(\mathbf{D}(\mathbf{M}))) = 0\)（对 \(i \neq 0\)），且同态 \(\alpha_{\mathrm{M}}\) 是双射的。

先取 M = A。映射 \(e : \Omega \to \mathbf{D}(A)\) 是一个拟同构（例 1），故映射 \(\mathbf{D}(e) : \mathbf{D}(\mathbf{D}(A)) \to \mathbf{D}(\Omega)\) 亦然。映射 \(\tilde{e} : A \to \mathbf{D}(\Omega)\) 是一个拟同构（例 2），且有 \(\mathbf{D}(e) \circ \alpha_{\mathbf{A}} = \tilde{c}\)；于是 \(\alpha_{A}\) 是拟同构，这就证明了此情形下的定理。由此推出，当 A-模 M 是有限生成的自由模时，\(\alpha_{M}\) 是拟同构。

转到一般情形；我们将对整数 n 作归纳，证明下列断言：

\((A_{n})\) 对每个有限生成的 A-模 M，同态 \(H_{i}(\alpha_{M})\) 当 \(i \leqslant n\) 时是双射的。

这也意味着 \(\mathrm{H}_{i}(\mathbf{D}(\mathbf{D}(\mathbf{M})))\) 在 \(i \neq 0\) 且 \(i \leqslant n\) 时为零，且 \(\alpha_{M}\) 当 \(n \geqslant 0\) 时是双射的。注意，\((\mathrm{A}_{n})\) 当 n \textless{} -d 时成立，其中 d 是复形 I 的长度：事实上，A-模 \(\mathbf{D}(\mathbf{D}(\mathbf{M}))_{i}\) 等于 \(\bigoplus_{p}\mathrm{Hom}_{\mathrm{A}}(\mathrm{Hom}_{\mathrm{A}}(\mathrm{M},\mathrm{I}^{p}),\mathrm{I}^{p-i})\)，故当 i \textless{} -d 及 i \textgreater{} d 时它为零。

我们来证明蕴含关系 \((\mathrm{A}_{n}) \Rightarrow (\mathrm{A}_{n+1})\)。设 M 为有限生成的 A-模。存在一个有限生成的自由 A-模 L 及一个正合列 \(0 \to N \xrightarrow{u} L \xrightarrow{v} M \to 0\)。序列 \(0 \to \mathbf{D}(M) \xrightarrow{\mathbf{D}(v)} \mathbf{D}(L) \xrightarrow{\mathbf{D}(u)} \mathbf{D}(N) \to 0\) 是正合的；同样，若令 \(u' = \mathbf{D}(\mathbf{D}(u))\) 与 \(v' = \mathbf{D}(\mathbf{D}(v))\)，则序列 \(0 \to \mathbf{D}(\mathbf{D}(N)) \xrightarrow{u'} \mathbf{D}(\mathbf{D}(L)) \xrightarrow{v'} \mathbf{D}(\mathbf{D}(M)) \to 0\) 是正合的。

由于 \(H_{i}(\mathbf{D}(\mathbf{D}(\mathrm{L})))\)（对 \(i \neq 0\)）为零，我们有同构

\[
\mathrm{H} _ {i} (\mathbf {D} (\mathbf {D} (\mathrm{M}))) \longrightarrow \mathrm{H} _ {i - 1} (\mathbf {D} (\mathbf {D} (\mathrm{N}))) \quad \text { pour } i \neq 0, 1;
\]

这就蕴含了蕴含关系 \((\mathrm{A}_{n})\Rightarrow(\mathrm{A}_{n+1})\)（对 \(n\neq-1\) 与 \(n\neq0\)）。考虑行正合的交换图表

\[
\begin{array}{c c c c c c c} 0 \to & N & \xrightarrow {u} & L & \xrightarrow {v} & M & \to 0 \\ & \alpha_ {N} \Bigg \downarrow & & \alpha_ {L} \Bigg \downarrow & & \alpha_ {M} \Bigg \downarrow \\ 0 \to H _ {1} (\mathbf {D} (\mathbf {D} (M))) & \xrightarrow {} H _ {0} (\mathbf {D} (\mathbf {D} (N))) & \xrightarrow {H _ {0} (u ^ {\prime})} & H _ {0} (\mathbf {D} (\mathbf {D} (L))) & \xrightarrow {H _ {0} (v ^ {\prime})} & H _ {0} (\mathbf {D} (\mathbf {D} (M))) \longrightarrow H _ {- 1} (\mathbf {D} (\mathbf {D} (N)) \end{array}
\]

其中 \(\alpha_{\mathrm{L}}\) 是双射的。若 \((\mathrm{A}_0)\) 成立，则同态 \(\alpha_{\mathrm{N}}\) 也是双射的，故 \(\mathrm{H}_0(u')\) 是单射的，且得到 \(\mathrm{H}_1(\mathbf{D}(\mathbf{D}(\mathbf{M}))) = 0\)，由此得 \((\mathrm{A}_1)\)。若 \((\mathrm{A}_{-1})\)

成立，则 \(\mathrm{H}_{-1}(\mathbf{D}(\mathbf{D}(\mathrm{N})))\) 为零，故 \(\mathrm{H}_{0}(v^{\prime})\) 是满射的，这蕴含 \(\alpha_{M}\) 是满射的。由于这对每个有限生成的 A-模 M 都成立，\(\alpha_{N}\) 也是满射的；由 I, § 1, 第 4 节，命题 2 的推论 2，\(\alpha_{M}\) 是双射的，于是 \((\mathrm{A}_{0})\) 成立。

于是 \((\mathrm{A}_{n})\) 对所有 n 都成立，这就证明了定理。

设 M 为有限生成的 A-模；令 \(c = \dim(A) - \dim_{A}(M)\)。令 \(\mathbf{D}'(M)\) 为 \(\mathbf{D}(M)\) 的下述子复形：它等于 \(\bigoplus_{i\in\mathcal{O}}\mathbf{D}(M)^{i}\bigoplus Z^{c}(\mathbf{D}(M))\)；并用

\[
j: \mathbf {D} ^ {\prime} (\mathrm{M}) \rightarrow \mathbf {D} (\mathrm{M})
\]

表示典范单射。由典范满射 \(\mathrm{Z}^{c}(\mathbf{D}(\mathrm{M}))\to\mathrm{H}^{c}(\mathbf{D}(\mathrm{M}))\) 与同构 \(\varphi(\mathrm{M},\mathrm{I})\) 导出复形的一个态射

\[
p: \mathbf {D} ^ {\prime} (\mathrm{M}) (- c) \rightarrow \operatorname{Ext} _ {\Lambda} ^ {c} (\mathrm{M}, \Omega);
\]

由于 \(\mathrm{H}^{i}(\mathbf{D}(\mathrm{M}))\) 当 i \textless{} c 时为零（第 1 节，命题 3 a)），\((\mathbf{D}^{\prime}(\mathrm{M})(-c), p)\) 是 \(\operatorname{Ext}_{\mathrm{A}}^{c}(\mathrm{M}, \Omega)\) 的一个左分解。由 A, X, 第 100 页，定理 1，我们有一个典范同构

\[
\varphi^ {0} \left(\mathbf {D} ^ {\prime} (\mathrm{M}) (- c), \mathrm{I}\right): \mathrm{H} _ {0} \left(\mathbf {D} \left(\mathbf {D} ^ {\prime} (\mathrm{M})\right)\right) \longrightarrow \operatorname{Ext} _ {\mathrm{A}} ^ {c} \left(\operatorname{Ext} _ {\mathrm{A}} ^ {c} (\mathrm{M}, \Omega), \Omega\right);
\]

把它与同态 \(\mathrm{H}_{0}(\mathbf{D}(j)) : \mathrm{H}_{0}(\mathbf{D}(\mathbf{D}(\mathrm{M}))) \to \mathrm{H}_{0}(\mathbf{D}(\mathbf{D}'(\mathrm{M})))\) 复合，便得到一个同态

\[
\mathrm{H} _ {0} (\mathbf {D} (\mathbf {D} (\mathrm{M}))) \longrightarrow \operatorname{Ext} _ {\mathrm{A}} ^ {c} \left(\operatorname{Ext} _ {\mathrm{A}} ^ {c} (\mathrm{M}, \Omega), \Omega\right),
\]

最后，通过与 \(\alpha_{\mathrm{M}}\) 复合，便得到一个典范同态

\[
\beta_ {\mathrm{M}}: \mathrm{M} \longrightarrow \operatorname{Ext} _ {\mathrm{A}} ^ {c} (\operatorname{Ext} _ {\mathrm{A}} ^ {c} (\mathrm{M}, \Omega), \Omega).
\]

推论.—— 若 A-模 M 是 Macaulay 模，则同态 \(\beta_{\mathrm{M}}\) 是双射的。

若 M 是 Macaulay 模，则 \(\Lambda\)-模 \(\mathrm{H}^i(\mathbf{D}(\mathrm{M}))\) 当 \(i \neq c\) 时为零（\(n^\circ 1\)，命题 3 的推论），故典范单射 \(j: \mathbf{D}'(\mathrm{M}) \to \mathbf{D}(\mathrm{M})\) 是拟同构；于是复形的态射 \(\mathbf{D}(j): \mathbf{D}(\mathbf{D}(\mathrm{M})) \to \mathbf{D}(\mathbf{D}'(\mathrm{M}))\) 是拟同构（\(\Lambda\), X, 第 86 页，命题 4）。从而 \(\mathrm{H}_0(\mathbf{D}(j))\) 是双射的；另一方面，由定理 1，\(\alpha_{\mathrm{M}}\) 是双射的，推论得证。

当 A-模 M 是有限长度时，A-模 \(\operatorname{Ext}_{\Lambda}^{c}(M,\Omega)\) 等同于 M 的 Matlis 对偶（参看 § 8, 第 5 节，例 3 及定理 3），由此重新得到 § 8, 第 4 节的命题 4。

